% Physical pages 3 and 4 of the concise study: the movement of the formulary,
% stated as an argument and not as a summary of what follows. The developed
% comparison of the three readings begins on the fifth page.

\section{The Propers: Themes and Movement}
\zlabel{triptych:concise:themes:start}

This Mass is about a call and what those who answer it bring with them. It
begins with God promising to hear whoever cries to him out of any tribulation,
reads the parable of a king whose invitation is refused and then carried to the
highways, and ends by asking that God's healing work free us from our
perversities and make us cleave to his commandments. At its centre is a man who
answered the call, came in, sat down at the king's table, and was found there
without a wedding garment.

The Introit gives God the first word. \latin{Salus populi ego sum, dicit
Dominus}: he is the salvation of the people, he will hear them \latin{de
quacumque tribulatione}, from whatever tribulation they cry to him, and he will
be their Lord \latin{in perpetuum}. Every clause of the antiphon is a promise,
and its words stand in no single verse of the Vulgate. The verse that follows
opens Ps~77, in which Asaph tells Israel its own history, ``That another
generation might know them'' and ``may not become like their fathers, a perverse
and exasperating generation''. The psalm's next verse, ``I will open my mouth in
parables'', is the one St Matthew cites as fulfilled in Jesus's speaking in
parables (Mt 13:35), and the Gospel of this Mass opens with Jesus speaking
\latin{in parabolis}. A people told the story of its fathers, so that it may
not become like them, is about to hear a parable about those first invited.

The Collect asks God to shut out \latin{universa nobis adversantia}, all that is
adverse, so that \latin{mente et corpore pariter expediti}, disencumbered in
mind and body alike, we may carry out \latin{quae tua sunt}, the things that are
his, \latin{liberis mentibus}, with free minds. Both halves of the petition return
at the end of the Mass: \latin{expediti} in the Postcommunion's
\latin{expediat}, and the things that are God's in its \latin{tuis \dots\
mandatis}, his commandments.

The Epistle begins in the middle of a sentence. St Paul has just told the
Ephesians to put off the old man (Eph 4:22), and the lesson completes the
thought: be renewed in the spirit of your mind, ``And put on the new man, who
according to God is created in justice and holiness of truth.'' Four commands
follow, and each puts off an act of the old man and puts on one of the new.
Lying gives way to truth spoken to the neighbour, ``For we are members one of
another.'' Anger is allowed its first motion and refused its duration: ``Let
not the sun go down upon your anger.'' The devil is given no place. And the
thief becomes a worker who gives: ``let him labour, working with his hands the
thing which is good, that he may have something to give to him that suffereth
need.'' The command is to clothe oneself. The verses after the lesson forbid
evil speech, warn ``grieve not the holy Spirit of God: whereby you are sealed
unto the day of redemption'', and end with forgiveness, ``even as God hath
forgiven you in Christ'' (4:29--32). The chapter from which the lesson is
cut opens with Paul beseeching his readers ``that you walk worthy of the
vocation in which you are called'' (4:1).

The Gradual is an evening prayer from Ps~140: ``Let my prayer be
directed as incense in thy sight; the lifting up of my hands, as evening
sacrifice.'' Its verb \latin{dirigatur}, may it be set straight, comes back in
the Communion's \latin{dirigantur}, and at a Mass celebrated with incense the
priest says the same verse again as he incenses the altar at the Offertory
(\work{Ordo Missae}, no.~1037). The Alleluia, the first verse of Ps~104, holds
three commands in a row: give glory, call upon his name, and ``declare his
deeds among the Gentiles''. The psalm goes on to recount Israel's history
from the covenant with Abraham to the gift of the land, and gives the reason for all of it in its last verse: ``That they might
observe his justifications, and seek after his law'' (v.~45). The Communion
will ask for exactly that.

The Gospel moves in two halves. In the first (Mt 22:2--10) a king twice sends
servants to those invited to his son's wedding. At first ``they would not
come''; then they make light of it and go off, ``one to his farm and another to his
merchandise'', and the rest seize the servants and kill them. The king destroys
the murderers and burns their city, declares the invited unworthy, and sends
his servants \latin{ad exitus viarum}, to the ends of the roads, to
call everyone they find; the hall is filled with ``both bad and good''. In the
second half (vv.~11--14) the king comes in to see those at table and finds one
man not clothed in a wedding garment. He calls him \latin{Amice}, friend, and
asks how he came in; \latin{At ille obmutuit}. The man is bound hand and foot
and cast into the exterior darkness, and the parable ends with a saying that
reaches beyond the story: ``For many are called, but few are chosen.'' The
Missal's opening names its hearers, \latin{principibus sacerdotum et
pharisaeis}, and so sets the parable in the temple at Jerusalem. The chief priests and ancients have
asked Jesus by what authority he acts (21:23); he has answered with the two sons
and the wicked husbandmen, and has said that ``the kingdom of God shall be taken
from you and shall be given to a nation yielding the fruits thereof'' (21:43);
and they knew that he spoke of them (21:45). The marriage feast is his third
answer, and after it the Pharisees go out to take counsel ``how to insnare him
in his speech'' (22:15).

The Offertory takes up the Introit's word \latin{tribulatio}, which occurs
nowhere else in the formulary: \latin{Si ambulavero in medio tribulationis,
vivificabis me}. The antiphon at the entrance promised a hearing from any
tribulation; the chant at the offering is sung by one who is still walking in
it. Where the Vulgate says that God has stretched forth his hand and that his
right hand has saved, the Missal sets both verbs in the future, so that the
whole verse looks forward. The Secret offers the gifts \latin{oculis tuae
maiestatis}, to the eyes of God's majesty, before which the Gradual asked the
prayer to rise \latin{in conspectu tuo}, and asks one thing, that they be
\latin{salutaria nobis}, saving to us. The image of offering before God's eyes
runs through this part of the Mass: at the offering of the chalice the priest
asks that it rise \latin{in conspectu divinae maiestatis tuae} (\work{Ordo
Missae}, no.~1032), and at the end of every Mass the
\latin{Placeat} asks in the same three words that the sacrifice offered
\latin{oculis tuae maiestatis} be acceptable to God (no.~1137).

The Communion sets a command beside a wish: \latin{Tu mandasti mandata tua
custodiri nimis: utinam dirigantur viae meae, ad custodiendas iustificationes
tuas}. The evening prayer asked to be set straight like incense; the Communion
asks that the ways of the one who prays be set straight towards what God
commands, and the justifications it names are those that Ps~104 gave as the
purpose of Israel's whole history. The Postcommunion names God's work by one
phrase, \latin{tua \dots\ medicinalis operatio}, thy healing working, and asks
two things of it. The first is the Collect's freedom again, now from something
within rather than from adversities without: \latin{a nostris perversitatibus
clementer expediat}. The second is the Communion's keeping again: \latin{tuis
semper faciat inhaerere mandatis}, that it make us cleave to his commandments.
The Mass that began with God promising to be his people's Lord \latin{in
perpetuum} ends with a prayer to hold to his commandments \latin{semper}.

Three words carry the whole: calling, hearing and keeping. The Lord hears
whoever cries to him, the Alleluia bids his people call upon his name, the king
calls twice and then calls everyone, and the Communion and Postcommunion ask to
keep and to cleave to what God commands. The Gospel's guest fails at the point where these meet. He was called, and
he came; he has nothing to say when the king speaks to him; and he is not
clothed for the feast to which he came.

Three readings of the whole Mass follow from three questions the Gospel raises,
each developed at length in the expansive study with four senses of its own.
The first asks what the guest who has come in must be wearing. St Gregory the
Great and St Augustine answer that the wedding garment is charity; St Jerome,
that it is the garment of the new man, made of the Lord's precepts kept, and
that the new man is Christ; and the reading finds that garment shown in the
Epistle's four acts and asked for in the Communion's wish. The second asks who
was called, by whom, and what became of the call when it was refused. It hears
the Mass from the Gospel's first half, the Introit's psalm and the Alleluia's
command, and St John Chrysostom, St
Irenaeus and St Hilary of Poitiers each read the parable as the call passing
from the first invited to the nations. The third asks where the feast is now
and how the assembly at this Mass takes its place at it. St Augustine and St
Gregory both say that this wedding is the present Church, with good and bad at
one table, and St John Chrysostom addresses its guests as those who have
partaken of the mysteries; the reading hears the Gradual, the Offertory, the
Secret and the Communion as the prayers of that assembly. The three answer
different questions and do not contradict one another, and the second supplies
the history in which the other two take place: the guests examined for the
garment and the assembly at the table are the ones gathered from the highways.
Within the readings the witnesses differ over what the garment is and where it
comes from, what the burned city was, where the highways lead, and
what the evening sacrifice is.
\zlabel{triptych:concise:themes:end}
